% =====================================================================
% FINAL SUBMISSION -- 31 October 2026.
% At most 8 pages excluding references, plus at most two appendix pages.
% Structured as a research paper. Every claim backed by prior results or by
% our own numbers.
% The challenges, limitations and future work section sits immediately before
% the references and is exempt from the page count, provided it contains only
% that and no results.
% Graded on the write-up, the code as a ZIP, and a viva with a presentation.
% Submit as ANLPDoomsday-Final.zip (this PDF + the code + the slides).
% =====================================================================
\documentclass[11pt]{article}
\usepackage{acl}
\usepackage{times}
\usepackage{latexsym}
\usepackage[T1]{fontenc}
\usepackage[utf8]{inputenc}
\usepackage{microtype}
\usepackage{amsmath,amssymb}
\usepackage{booktabs}

\newcommand{\idf}{\mathrm{idf}}
\newcommand{\sfreq}{\mathrm{sf}}

\title{Predicting When Collection Statistics Help Top-$k$ Selection:\\
From Retrieval to Attention and Expert Routing}

\author{
  \textbf{Team ANLP Doomsday} \\[2pt]
  Aviral Gupta \quad Mohit Kumar Singh \quad Anirudh Sankar \quad Arihant Tripathy \\
  \texttt{\{aviral.gupta, mohit.singh, anirudh.sankar, arihant.tripathy\}@research.iiit.ac.in} \\
  IIIT Hyderabad, India
}

\begin{document}
\maketitle

\begin{abstract}
% Written last. Check it against the body on the direction of every reported
% effect, which is a mistake that has been made before.
\end{abstract}

\section{Introduction}
\label{sec:intro}

\section{Motivation and Research Gap}
\label{sec:motivation}
% The shortcomings of existing work that suggest this direction. Counted in the page
% limit.

\section{Contributions}
\label{sec:contributions}
% How the project addresses those shortcomings, against measured results.

\section{Background}
\label{sec:background}

\section{Related Work}
\label{sec:related}

\section{Method}
\label{sec:method}
% The modular scorer and the four switches, the eight bases at two
% granularities, the six diagnostics, and the predictor.

\section{Experimental Setup}
\label{sec:setup}
% Datasets, budgets, hyperparameters, and the pre-registered analysis: primary
% comparison all-on against all-off per representation, component ablations
% secondary, paired bootstrap over queries with Holm-Bonferroni within each
% family, effect sizes alongside p-values.

\section{Results}
\label{sec:results}
% Setting A grid, the diagnostics-to-margin regression under
% leave-one-basis-family-out validation, then the transfer to Settings B and C
% with no refitting.

\section{Analysis}
\label{sec:analysis}
% Which component of the scorer carries the effect, where the transfer holds and
% where it inverts, and what that says about frequency against importance inside
% a model. The collapse condition goes here.

\section{Conclusion}
\label{sec:conclusion}

% Exempt from the page count. Nothing but challenges, limitations and future
% work; a result placed here would be read as a limitation.
\section*{Challenges, Limitations and Future Work}
\label{sec:limitations}

%\bibliography{refs}

\appendix
% Two pages at most.

\end{document}
